% ================== Chapitre 3 : isolation et confidentialite ===============
\chapter{Isolation et confidentialité des noms}
\label{ch:isolation}

\section{Expression du besoin et scénarios à empêcher}

L'exigence \exig{net-fr-008} porte sur la confidentialité des noms entre
domaines. Le document 03 la précise sans ambiguïté : le mécanisme de résolution
doit empêcher la lecture d'une zone d'infrastructure ou d'un autre tenant par
les clients non autorisés, et cacher un enregistrement dans le portail ne suffit
pas. Le test \opt{S04} formule la vérification attendue : depuis une machine du
VNet USER, une requête vers la zone d'infrastructure et vers la zone d'un tenant
voisin doit se solder par un refus ou une absence d'information.

Trois scénarios doivent être couverts, et chacun appelle une contre-mesure
distincte.

\begin{table}[H]
\centering
\small
\caption{Scénarios de divulgation et contre-mesures}\label{tab:scenarios}
\begin{tabular}{@{}R{0.30}R{0.33}R{0.37}@{}}
\hautfilet
\textbf{Scénario} & \textbf{Ce qui serait divulgué} &
\textbf{Contre-mesure retenue}\\
\medfilet
Interrogation directe d'un serveur autoritatif par une machine tenant &
  Le serveur autoritatif répondrait pour toute zone hébergée, y compris celles
  des autres tenants. &
  Filtrage réseau : les adresses \ip{10.30.4.10} et \ip{10.30.4.11} ne sont pas
  joignables depuis un VNet tenant (section~\ref{sec:flux}).\\
Transfert de zone demandé par un client quelconque &
  Un transfert divulgue l'intégralité d'une zone en une requête. &
  Transferts restreints à ns2 par adresse et signés par une clé TSIG
  (section~\ref{sec:zones}).\\
Requête vers une zone étrangère au travers du résolveur partagé &
  Le résolveur résout toutes les zones internes pour tous ses clients. &
  Politique de sélection des zones indexée sur l'adresse source
  (section~\ref{sec:vue}).\\
\basfilet
\end{tabular}
\end{table}

Une énumération par requêtes successives reste théoriquement possible à
l'intérieur de la vue d'un tenant, c'est-à-dire sur ses propres noms. Ce résidu
est accepté : il ne franchit aucune frontière de tenant.

\section{Politique de sélection des zones par adresse source}
\label{sec:vue}

Les deux options ouvertes par le document 03 sont de séparer physiquement les
points de résolution par tenant, ou d'appliquer une politique démontrée de
sélection des zones. La première option supposerait cinquante paires de
résolveurs, ce que le matériel disponible exclut. La seconde est donc retenue.

Son principe est direct. La relation \eqref{eq:prefixe} associe le tenant $n$ au
préfixe \ip{10.30.(63+n).0/24}. L'adresse source d'une requête détermine donc
sans ambiguïté le tenant émetteur, et par conséquent les deux seules zones
internes qu'il est autorisé à résoudre. En notant $\mathcal{A}$ la réunion des
préfixes d'administration et $o_3(a)$ le troisième octet de l'adresse $a$,
l'étiquette de vue s'écrit
\begin{equation}
  \tau(a) =
  \begin{cases}
    1 & \text{si } a \in \mathcal{A},\\[3pt]
    1000 + \bigl(o_3(a) - 63\bigr)
      & \text{si } a \in \texttt{10.30.64.0/18}
        \text{ et } 64 \leqslant o_3(a) \leqslant 113,\\[3pt]
    999 & \text{sinon,}
  \end{cases}
  \label{eq:tau}
\end{equation}
et l'ensemble des zones internes résolubles par un client d'étiquette
$\tau$ vaut
\begin{equation}
  \mathcal{Z}(\tau) =
  \begin{cases}
    \text{toutes les zones internes} & \text{si } \tau = 1,\\[3pt]
    \bigl\{\, \mathcal{D}(\tau-1000),\ \mathcal{R}(\tau-1000) \,\bigr\}
      & \text{si } \tau > 1000,\\[3pt]
    \varnothing & \text{si } \tau = 999 .
  \end{cases}
  \label{eq:zones}
\end{equation}

Toute demande portant sur l'espace interne et sortant de $\mathcal{Z}(\tau)$
reçoit un code \opt{REFUSED}, qui ne révèle pas l'existence du nom demandé,
contrairement à une réponse \opt{NXDOMAIN}.

\begin{encadredecision}
\titreencadredec{Point de conception critique : où placer le contrôle}
Une vérification placée dans le seul point d'entrée de résolution serait
insuffisante. Le résolveur consulte son cache paquet avant d'engager la
résolution : une réponse construite pour un tenant pourrait être restituée telle
quelle à un autre, et la politique serait contournée par le cache lui-même.

\smallskip
La fonction \opt{gettag()} est évaluée avant toute consultation de cache et sa
valeur de retour entre dans la clé du cache paquet. Le contrôle est donc
construit en deux temps : \opt{gettag()} calcule l'étiquette $\tau(a)$ de
l'équation \eqref{eq:tau}, puis \opt{preresolve()} applique la règle
\eqref{eq:zones} en s'appuyant sur cette étiquette. Deux vues distinctes ne
peuvent alors jamais partager une entrée de cache.
\end{encadredecision}

\begin{figure}[H]
\centering
\figChemin
\caption{Chemin d'une requête issue d'une machine tenant et points
  d'application de la politique}
\label{fig:chemin}
\end{figure}

\begin{table}[H]
\centering
\small
\caption{Vues définies par la politique de sélection des zones}\label{tab:vues}
\begin{tabular}{@{}C{0.12}R{0.41}R{0.47}@{}}
\hautfilet
\textbf{$\tau$} & \textbf{Adresses source} & \textbf{Portée de résolution}\\
\medfilet
1 & Plan de gestion \ip{10.30.0.0/24}, VNet Services \ip{10.30.4.0/24},
  Monitoring \ip{10.30.5.0/24}, pool VPN administrateur \ip{10.30.16.32/28} &
  Vue administration : résolution de l'ensemble de l'espace interne et de
  l'espace public.\\
$1000+n$ & Préfixe du tenant $n$, soit \ip{10.30.(63+n).0/24} &
  Zone directe $\mathcal{D}(n)$, zone inverse $\mathcal{R}(n)$, et espace
  public.\\
999 & Toute autre source acceptée par \opt{allow-from}, notamment le pool VPN
  utilisateur \ip{10.30.16.16/28} &
  Espace public uniquement ; tout nom interne reçoit \opt{REFUSED}.\\
aucune & Source hors \opt{allow-from} &
  La requête n'est pas traitée par le résolveur.\\
\basfilet
\end{tabular}
\end{table}

\section{Filtrage réseau et matrice de flux}
\label{sec:flux}

La politique applicative ne remplace pas le filtrage. Les règles suivantes
complètent les politiques de référence du document 02 et sont appliquées sur les
passerelles pfSense des VNets concernés, conformément aux interfaces IF-INT-05
et IF-INT-07.

\begin{longtable}{@{}R{0.21}R{0.17}R{0.13}R{0.16}R{0.33}@{}}
\caption{Matrice de flux du service DNS}\label{tab:matrice}\\
\hautfilet
\textbf{Source} & \textbf{Destination} & \textbf{Service} &
\textbf{Décision} & \textbf{Motif}\\
\medfilet
\endfirsthead
\hautfilet
\textbf{Source} & \textbf{Destination} & \textbf{Service} &
\textbf{Décision} & \textbf{Motif}\\
\medfilet
\endhead
\basfilet
\endfoot
VNet tenant \ip{10.30.(63+n).0/24} & \ip{10.30.4.12}, \ip{10.30.4.13} &
  UDP et TCP 53 & Autoriser & Seul point de résolution offert aux tenants\\
VNet tenant & \ip{10.30.4.10}, \ip{10.30.4.11} & tout &
  Refuser et journaliser & Les autoritatifs ne sont pas exposés aux tenants\\
VNet tenant & \ip{10.30.0.0/24} & tout & Refuser &
  Règle existante du document 02, rappelée ici\\
\ip{10.30.4.12}, \ip{10.30.4.13} & \ip{10.30.4.10}, \ip{10.30.4.11} &
  UDP et TCP 53 & Autoriser & Redirection de zone interne\\
\ip{10.30.4.12}, \ip{10.30.4.13} & Internet & UDP et TCP 53 &
  Autoriser via PF-WAN & Récursion sur l'espace public, journalisée\\
\ip{10.30.4.11} (ns2) & \ip{10.30.4.10} (ns1) & TCP 53 & Autoriser &
  Transfert de zone signé TSIG\\
\ip{10.30.4.10} (ns1) & \ip{10.30.4.11} (ns2) & UDP 53 & Autoriser &
  Message NOTIFY\\
\ip{10.30.4.41} (agent) & \ip{10.30.4.10} & TCP 8443 & Autoriser &
  API autoritative derrière terminaison TLS\\
\ip{10.30.4.41} (agent) & ns1, ns2, rec1, rec2 & UDP et TCP 53 & Autoriser &
  Contrôle de publication sur les quatre instances\\
\ip{10.30.0.0/24}, \ip{10.30.16.32/28} & \ip{10.30.4.12}, \ip{10.30.4.13} &
  UDP et TCP 53 & Autoriser & Vue administration\\
\ip{10.30.16.16/28} (VPN utilisateur) & \ip{10.30.4.12}, \ip{10.30.4.13} &
  UDP et TCP 53 & Autoriser &
  Admis par \opt{allow-from}, mais limité à l'espace public par la vue 999
  (section~\ref{sec:vpn})\\
\ip{10.30.5.10} (Monitoring) & ns1, ns2, rec1, rec2 & TCP 8443 & Autoriser &
  Lecture des métriques, initiée par Monitoring\\
ns1, ns2, rec1, rec2 & \ip{10.30.5.0/24} & tout & Refuser &
  Aucune initiation vers Monitoring, conformément à IF-NS-08\\
Toute autre source & ns1, ns2, rec1, rec2 & tout & Refuser &
  Refus par défaut\\
\end{longtable}

Deux points méritent attention. Le refus du trafic des tenants vers les
autoritatifs doit être journalisé : une hausse de ce compteur signale soit une
erreur de distribution des résolveurs, soit une tentative délibérée. Par
ailleurs, le port 53 doit être ouvert en TCP autant qu'en UDP, faute de quoi
toute réponse tronquée deviendrait irrécupérable, avec un symptôme limité à
certaines requêtes seulement.

\section{Cas particulier des profils VPN}
\label{sec:vpn}

L'exigence \exig{net-fr-017} demande que le DNS soit accessible selon le profil
VPN. Le plan d'adressage distingue trois pools : administrateur en
\ip{10.30.16.32/28}, utilisateur en \ip{10.30.16.16/28} et site à site en
\ip{10.30.16.48/28}, les deux derniers étant classés en extension.

Le pool administrateur ne pose pas de difficulté : il est traité comme le plan
de gestion et reçoit la vue administration, sous réserve de l'autorisation
délivrée par IAM au moment de l'ouverture de session.

Le pool utilisateur pose en revanche un problème de conception qu'il faut
énoncer plutôt que masquer. Une politique fondée sur l'adresse source ne peut
pas distinguer deux utilisateurs de tenants différents qui reçoivent leur
adresse d'un même pool partagé de quatorze adresses. Trois réponses sont
possibles.

\begin{enumerate}
  \item Attribuer une sous-plage distincte par tenant, ce que la taille du pool
    actuel ne permet pas pour cinquante tenants.
  \item Faire porter l'identité du tenant par le concentrateur VPN, qui
    attribuerait une adresse stable par session et transmettrait la
    correspondance au résolveur, au titre de l'interface IF-EXT-07.
  \item Refuser toute résolution interne au pool utilisateur et médier les accès
    par le Backend, conformément à IF-EXT-06.
\end{enumerate}

La troisième réponse est retenue pour le périmètre minimal, parce qu'elle est la
seule applicable sans donnée supplémentaire et parce qu'elle est cohérente avec
le statut d'extension du pool utilisateur. L'étiquette $\tau = 999$ de
l'équation \eqref{eq:tau} matérialise ce refus. La deuxième réponse constitue
l'évolution naturelle dès que la correspondance session vers tenant sera
fournie.
